\documentclass{article}
\usepackage{amsmath, amssymb, amsthm}

\setlength{\parindent}{0pt}

\newtheorem{claim}{Claim}

\begin{document}

\begin{claim}
    A busy student must complete $3$ problem sets before doing laundry. Each problem set 
    requires $1$ day with probability $\frac{2}{3}$ and $2$ days with probability 
    $\frac{1}{3}$. Let $B$ be the number of days a busy student delays laundry. Then 
    $\mathrm{E}[B] = 4$.
\end{claim}

\begin{proof}
    Let $X_1$, $X_2$, and $X_3$ denote the number of days required to complete the first, 
    second, and third problem sets respectively. Then 
    \[
    B = X_1 + X_2 + X_3
    \]
    By linearity of expectation,
    \[
    \mathrm{E}[B] = \mathrm{E}[X_1 + X_2 + X_3] 
    = \mathrm{E}[X_1] + \mathrm{E}[X_2] + \mathrm{E}[X_3]
    \]
    Since $X_1$, $X_2$, and $X_3$ have the same probability distribution, 
    \[
    \mathrm{E}[X_1] = \mathrm{E}[X_2] = \mathrm{E}[X_3]
    \]
    Hence it suffices to compute 
    $\mathrm{E}[X_1]$. By the definition of expectation,
    \[
    \mathrm{E}[X_1] = 1 \cdot \Pr\{X_1 = 1\} + 2 \cdot \Pr\{X_1 = 2\}
    \]
    Plugging $\Pr\{X_1 = 1\} = \frac{2}{3}$, $\Pr\{X_1 = 2\} = \frac{1}{3}$ into the above 
    equation gives
    \[
    \mathrm{E}[X_1] = 1 \cdot \frac{2}{3} + 2 \cdot \frac{1}{3} = \frac{4}{3}
    \]
    Therefore,
    \[
    \mathrm{E}[B] = \mathrm{E}[X_1 + X_2 + X_3] 
    = \mathrm{E}[X_1] + \mathrm{E}[X_2] + \mathrm{E}[X_3] 
    = \frac{4}{3} + \frac{4}{3} + \frac{4}{3} = 4
    \]

    \medskip

    Hence 
    \[
    \mathrm{E}[B] = 4
    \]
\end{proof}

\begin{claim}
    A relaxed student rolls a fair, $6$-sided die in the morning. If he rolls a $1$, then he 
    does his laundry immediately (with zero days of delay). Otherwise, he delays for one 
    day and repeats the experiment the following morning. Let $R$ be the number of days a 
    relaxed student delays laundry. Then $\mathrm{E}[R] = 5$.
\end{claim}

\begin{proof}
    Let $D$ denote the outcome of this morning's die roll. We condition on the result of this 
    morning's roll. If $D = 1$, the student does not delay his laundry, contributing $0$ days. 
    If $D \neq 1$, the student delays for one day and the remaining process has the same 
    distribution as the original process, since the next morning's roll is independent of all 
    previous rolls. Thus, the expected delay from that point onward is again $\mathrm{E}[R]$. 
    Therefore, by conditioning on the first roll,
    \[
    \mathrm{E}[R] = 0 \cdot \Pr\{D = 1\} + (1 + \mathrm{E}[R]) \cdot \Pr\{D \neq 1\}
    \]
    Since the die is fair, there is a probability of $\frac{1}{6}$ of rolling a $1$ on the 
    first try, and a probability of $\frac{5}{6}$ otherwise. Plugging into the above equation 
    gives
    \[
    \mathrm{E}[R] = 0 \cdot \frac{1}{6} + (1 + \mathrm{E}[R]) \cdot \frac{5}{6}
    \]
    Solving for $\mathrm{E}[R]$ gives
    \[
    \begin{aligned}
    \mathrm{E}[R] &= 0 \cdot \frac{1}{6} + (1 + \mathrm{E}[R]) \cdot \frac{5}{6} \\
    \mathrm{E}[R] &= \frac{5}{6} + \frac{5}{6}\mathrm{E}[R] \\
    \frac{1}{6}\mathrm{E}[R] &= \frac{5}{6} \\
    \mathrm{E}[R] &= 5
    \end{aligned}
    \]

    \medskip

    Hence
    \[
    \mathrm{E}[R] = 5
    \]
\end{proof}

\begin{claim}
    Before doing laundry, an unlucky student must recover from illness for a number of days 
    equal to the product of the numbers rolled on two fair, $6$-sided dice. Let $U$ be the
    expected number of days an unlucky student delays laundry. Then $\mathrm{E}[U] = \frac{49}{4}$.
\end{claim}

\begin{proof}
    Let $D_1$ and $D_2$ be the outcomes of the first and second die rolls respectively. 
    Then
    \[
    U = D_1D_2
    \]
    By the definition of expectation,
    \[
    \mathrm{E}[U] = \mathrm{E}[D_1D_2]
    \]
    Because $D_1$ and $D_2$ are independent,
    \[
    \mathrm{E}[D_1D_2] = \mathrm{E}[D_1]\mathrm{E}[D_2]
    \]
    Since $D_1$ and $D_2$ have the same probability distribution,
    \[
    \mathrm{E}[D_1] = \mathrm{E}[D_2]
    \]
    Hence it suffices to compute $\mathrm{E}[D_1]$. By the definition of expectation, 
    \[
    \mathrm{E}[D_1] = \sum_{x=1}^{6} x \cdot \Pr\{D_1 = x\}
    \]
    Since the die is fair,
    \[
    \Pr\{D_1 = x\} = \frac{1}{6}
    \]
    Thus
    \[
    \mathrm{E}[D_1] = \sum_{x=1}^{6} x \cdot \frac{1}{6} = \frac{7}{2}
    \]
    Therefore
    \[
    \mathrm{E}[U] = \mathrm{E}[D_1D_2] = \mathrm{E}[D_1]\mathrm{E}[D_2] 
    = \frac{7}{2} \cdot \frac{7}{2} = \frac{49}{4}
    \]

    \medskip

    Hence
    \[
    \mathrm{E}[U] = \frac{49}{4}
    \]
\end{proof}

\begin{claim}
    A student is busy with probability $\frac{1}{2}$, relaxed with probability 
    $\frac{1}{3}$, and unlucky with probability $\frac{1}{6}$. Let $D$ be the number 
    of days the student delays laundry. Then $\mathrm{E}[D] = \frac{137}{24}$.
\end{claim}

\begin{proof}
    Let $T$ denote the student's type, where $T = B$, $T = R$, and $T = U$ denote 
    that the student is busy, relaxed, and unlucky respectively. Then by the law of 
    total expectation, 
    \[
    \begin{aligned}
    \mathrm{E}[D] &= \mathrm{E}[D \mid T = B]\Pr\{T = B\} \\
    &\quad + \mathrm{E}[D \mid T = R]\Pr\{T = R\} \\
    &\quad + \mathrm{E}[D \mid T = U]\Pr\{T = U\}
    \end{aligned}
    \]
    By the previous claims,
    \[
    \begin{aligned}
    \mathrm{E}[D \mid T = B] &= 4 \\
    \mathrm{E}[D \mid T = R] &= 5 \\
    \mathrm{E}[D \mid T = U] &= \frac{49}{4}
    \end{aligned}
    \]
    and $\Pr\{T = B\} = \frac{1}{2}$, $\Pr\{T = R\} = \frac{1}{3}$, and 
    $\Pr\{T = U\} = \frac{1}{6}$. Substituting these values into the above equation 
    gives
    \[
    \mathrm{E}[D] = 4 \cdot \frac{1}{2} + 5 \cdot \frac{1}{3} + 
    \frac{49}{4} \cdot \frac{1}{6} = \frac{137}{24}
    \]

    \medskip

    Hence
    \[
    \mathrm{E}[D] = \frac{137}{24}
    \]
\end{proof}

\end{document}
